\section{Mechanism: The Interaction-Zone Characterization}
\label{sec:mechanism}

Section~\ref{sec:results} documented four facts: the moderation concentrates in mid-IO ($\hat\gamma_{\text{mid}} = -0.0256$, $t = -2.24$); it relocates to low-IO when the bottom 50 percent of firms by market cap are dropped ($\hat\gamma_{\text{IO1}} = -0.0441$, $t = -2.04$); the per-state coefficient tracks state literacy in the cell that carries the moderation ($r = -0.28$ full-panel mid-IO, $r = -0.29$ above-median IO1) but not in high-IO ($r = -0.07$); and the within-mid-IO IV-decile decomposition fails to reject decile equality ($\chi^2\,p = 0.94$), with decile 10 $\approx 0$. This section develops the interaction-zone characterization those facts identify and discusses what the design does and does not separate.

\subsection{The interaction-zone reading}
\label{sec:zone}

The moderation lives where two conditions hold at once: retail demand large enough to move prices, and institutional activity sufficient to translate that demand into cross-state return differences. Where neither holds (micro-caps held mostly by retail), mispricing accumulates but does not load on cross-state literacy variation; where institutions dominate (large high-IO firms), arbitrage absorbs retail demand and the marginal investor is not a local trader. The interaction zone --- mid-cap firms with mixed ownership --- is where both hold. Its IO label is sample-conditional: under the full-panel size composition it falls in mid-IO (mid-IO firms being predominantly mid-cap), and under the above-median composition it migrates to low-IO (``low IO'' among large firms meaning retail-tilted at substantial size). Table~\ref{tab:migration} traces this cell by cell; the high-IO null under both compositions is the predicted-null check the reading passes, and the IV-decile null within mid-IO (Section~\ref{sec:mech_iv_decile}) is consistent, since a size-and-ownership channel predicts no IV-monotonicity.

\subsection{What the IV-decile null rules out}
\label{sec:mech_iv_decile}

The canonical microfoundation in the limits-to-arbitrage literature \citep{shleiferVishny1997, pontiff2006} predicts that arbitrage costs scale with idiosyncratic risk, so literacy-modulated retail demand should displace more mispricing where arbitrage is more expensive. Within mid-IO, the within-month NYSE IV-decile decomposition tests this prediction directly (full decile table in IA~\ref{ia:ivwithinmidio}), and rejects it. Decile 10 --- the highest-IV decile, where the canonical microfoundation predicts the strongest effect --- returns $\hat\gamma = -0.035$ ($t = -1.27$), not significant and far from monotone, and a joint equality test across deciles fails to reject ($\chi^2(9)$, $p = 0.94$). The largest negative coefficient sits in decile 7, and the across-decile pattern is non-monotone; under the joint specification the Spearman $\rho(\text{decile}, \hat\gamma_d) = -0.67$ ($p = 0.033$) is a directional decline with IV, but with no concentration at the highest IV decile. The IV-decile margin is informative for what it rules out (canonical limits-to-arbitrage) but does not positively identify the interaction-zone reading; that identification comes from the size-restriction migration of Section~\ref{sec:migration}.

\subsection{High-IO null and college horse-race}
\label{sec:mech_hi_io}

The high-IO cross-state null ($\hat\beta_{\text{IO3}} = +0.018$, $t = +1.22$, $r = -0.07$; Holm $p_{\text{FWER}} = 0.996$) is the second confirmatory pattern the reading anticipates: high-IO firms are held by nationally diversified institutions indifferent to headquarters state, so state-level literacy variation says nothing about the marginal investor, and whatever retail demand exists is arbitraged away. Two diagnostics rule out the obvious alternatives. Literacy is not a proxy for education --- a state college-attainment horse-race within mid-IO leaves literacy sharp ($\hat\gamma_{\text{lit}} = -0.0350$, $t = -2.97$) and college null ($+0.00076$, $t = +0.65$) despite a $\rho = 0.62$ state-level correlation --- and the mid-IO concentration is not a variance artifact (mid-IO has lower identifying variation than low-IO yet a larger slope; Section~\ref{sec:robustness}).

\subsection{What the design does not separate}
\label{sec:mech_what_not}

The interaction-zone characterization is identified through the size-restriction migration and the cross-state gradient, but it does not separate three within-family microfoundations that share the same size-ownership conditioning: literacy-modulated retail demand \emph{intensity} (the reading the paper most naturally adopts, with the cross-state gradient direction consistent with it), literacy-modulated retail \emph{attention} (less momentum-following demand pressure in high-literacy states), and literacy-modulated retail \emph{timing} (a clientele-timing channel generating the same gradient). The three make different predictions for retail trade-flow data (volume, turnover, cross-spectral signatures) not available in the current panel; the TAQ-based Boehmer-Jones-Zhang retail trade indicator would separate them, a natural next-paper extension.

\subsection{Population-scope limitation from the PIT correction}
\label{sec:mech_population}

The characterization holds on the stable-HQ subsample, on which the local-investor unit of analysis is held fixed. Under the point-in-time 10-K-header reassignment, the aggregate three-way collapses roughly thirty-nine fold ($-0.0117 \to -0.0003$) and the mid-IO coefficient twenty-six fold ($-0.0321 \to -0.00121$): the directional within-tercile pattern survives but the magnitudes do not. The interaction-zone characterization is therefore a statement about stable-headquarters firms, not the PIT-corrected full panel. Section~\ref{sec:robustness} reports the full PIT diagnostic.
